\chapter{Related Work}
\label{chap:related}
% FIIT: literature review and comparison of existing approaches; 1--1.5 pages.

\section{Side-Channel Analysis of Neural Networks}

Physical side channels reveal the architecture, inputs and parameters of a
neural network~\cite{horvath2024sok}. Batina et al.~\cite{batina2019csi}
recovered both architecture and weights, but limited weight candidates to a
precision of $10^{-2}$. Joud et al.~\cite{joud2022practical} and Nozaki and
Kobara~\cite{nozaki2024power} raised the precision with iterative CEMA and
IEEE-754-aware correlation attacks, yet Joud et al.\ still report small errors
rather than exact values. Exact extraction exists under stronger assumptions,
using subnormal timing~\cite{gongye2020floating} or injected
faults~\cite{breier2021sniff}.

\section{Template Attacks}

A template attack first characterizes the leakage of an identical device and
then uses this model against the target~\cite{chari2002template}. Given a
leakage class $c$, the leakage vector $\mathbf{t}$ at the points of interest is
modeled as
\begin{equation}
	\mathbf{t} \mid c \sim
	\mathcal{N}(\boldsymbol{\mu}_c,\boldsymbol{\Sigma}_c),
\end{equation}
and candidates are ranked by likelihood, using several samples and their noise
correlation jointly. Later work replaces the Gaussian model with machine-learning
models such as two-dimensional CNNs~\cite{tu2021template}. In our implementation, the
covariance matrix is pooled across leakage classes. Wei et
al.~\cite{wei2018iknow} used a ``power template'' to recover inputs, but it is
a lookup from power features to pixels. To the best of our knowledge, Gaussian
templates have not been applied to recovering \texttt{float32} weights.
